%==============================================================================
% Quality Ladders & Schumpeterian Growth Cheatsheet
% Lecture 5 — Macroeconomics I (BSE)
% Exam-Focused Reference
%==============================================================================

\documentclass[a4paper,11pt]{article}

\usepackage{geometry}
\geometry{margin=1in, headheight=15pt}

\usepackage{amsmath}
\usepackage{amssymb}
\usepackage{mathtools}
\usepackage{graphicx}

\usepackage{xcolor}
\definecolor{blue3}{RGB}{0,0,180}
\definecolor{green3}{RGB}{0,120,0}
\definecolor{orange3}{RGB}{200,100,0}

\usepackage{siunitx}
\sisetup{detect-all}

\usepackage[colorlinks=true,linkcolor=blue3,citecolor=blue3,urlcolor=blue3]{hyperref}

\usepackage{cleveref}

\usepackage{tcolorbox}
\tcbuselibrary{skins,breakable}

\newtcolorbox{derivationbox}[1][]{
    colback=blue!10,
    colframe=blue3,
    coltitle=blue3,
    title={Derivation:},
    fonttitle=\bfseries,
    #1
}

\newtcolorbox{resultbox}[1][]{
    colback=green!10,
    colframe=green3,
    coltitle=green3,
    title={Result Summary:},
    fonttitle=\bfseries,
    #1
}

\newtcolorbox{warningbox}[1][]{
    colback=orange!10,
    colframe=orange3,
    coltitle=orange3,
    title={Important:},
    fonttitle=\bfseries,
    #1
}

\usepackage{booktabs}
\usepackage{enumitem}

% Custom notation commands
\newcommand{\pibar}{\overline{\pi}}
\newcommand{\Xtilde}{\tilde{X}}

\title{\textbf{Quality Ladders \& Schumpeterian Growth}\\[0.3em]
       \large Lecture 5 Cheatsheet --- Macroeconomics I (BSE)}
\author{Exam Reference}
\date{}

\begin{document}

\maketitle
\tableofcontents
\clearpage
\pagestyle{plain}

%==============================================================================
\section{Notation and Definitions}
%==============================================================================

\subsection*{Core Parameters and Variables}
\begin{align*}
& N = \text{number of intermediate goods (fixed)} & q &> 1 = \text{quality step size}\\
& k(s) = \text{quality rung of good } s & \alpha &\in (0,1) = \text{intermediate input share}\\
& L = \text{total labor (inelastic)} & L_Y &= \text{labor in final good sector}\\
& r = \text{interest rate} & \rho &= \text{household discount rate}\\
& \sigma = \text{intertemporal elasticity of substitution} & p(s) &= \text{innovation (arrival) rate}\\
& \xi = \text{R\&D cost parameter} & Z(s) &= \text{R\&D resources for good }s\\
& w = \text{wage} & \pi &= \text{monopoly profits}\\
& V(k) = \text{present value of innovation} & Q &= \text{aggregate quality index}\\
& \gamma = g_Y = \text{BGP growth rate}
\end{align*}

\subsection*{Key Derived Objects}
\begin{align*}
\Xtilde_t(s) &:= q^{k(s)} X_t(s) \quad \text{(quality-adjusted intermediate input)}\\
Q_t &:= \int_0^N q^{k(s)\frac{\alpha}{1-\alpha}} ds \quad \text{(aggregate quality index)}\\
\pibar &:= \frac{1-\alpha}{\alpha}\,\alpha^{\frac{2}{1-\alpha}} L \quad \text{(normalized per-sector profit)}
\end{align*}

\subsection*{Key Assumption: Innovation Step Size}
\begin{tcolorbox}[colback=orange!10,colframe=orange3,title=Critical Assumption]
The model assumes $q \geq \dfrac{1}{\alpha}$. Under this assumption, the new innovator can always
profitably undercut the old incumbent and become the \emph{sole} monopolist. This justifies
full creative destruction (no simultaneous coexistence of old and new producer).
\end{tcolorbox}

\subsection*{Differences from Romer (Expanding Varieties)}
\begin{center}
\begin{tabular}{lll}
\toprule
\textbf{Feature} & \textbf{Romer (Varieties)} & \textbf{Aghion-Howitt (Quality)} \\
\midrule
Growth mechanism & New goods ($N$ grows) & Quality improvements ($q$ rises) \\
$N$ (varieties) & Grows over time & Fixed \\
Incumbent fate & Coexists with entrants & Destroyed by entrant \\
Business stealing & No & Yes \\
Creative destruction & No & Yes \\
Externality sign & Positive only & Positive + Negative \\
Net welfare bias & Under-innovation & Ambiguous \\
\bottomrule
\end{tabular}
\end{center}

\clearpage

%==============================================================================
\section{Model Setup: Final Good and Quality-Adjusted Intermediates}
%==============================================================================

\subsection*{Final Good Production}

The representative final-good firm solves:
\[
\max_{L,\, X_t(s)} \; Y_t - w_t L - \int_0^N P_t(s)\, X_t(s)\, ds
\]
subject to the technology:
\begin{align*}
Y_t = L^{1-\alpha} \int_0^N \bigl(q^{k(s)} X_t(s)\bigr)^\alpha ds
\tag{Final Good}
\end{align*}

where $N$ is the (fixed) measure of intermediate goods and $q^{k(s)}$ embeds quality into productivity.

\subsection*{First-Order Conditions of Final Good Firm}

\textbf{Labor demand:}
\begin{align*}
w_t = (1-\alpha)\,L^{-\alpha}\int_0^N \Xtilde_t^\alpha(s)\,ds = (1-\alpha)\,\frac{Y_t}{L}
\tag{Labor FOC}
\end{align*}

\textbf{Intermediate input demand (inverse demand for $X_t(s)$):}
\begin{align*}
P_t(s) = \alpha\, L^{1-\alpha}\, q^{\alpha k(s)}\, X_t(s)^{\alpha-1}
\tag{Demand}
\end{align*}

Inverting the demand function:
\begin{align*}
\boxed{X_t(s) = L\left(\frac{\alpha\, q^{\alpha k(s)}}{P_t(s)}\right)^{\frac{1}{1-\alpha}}}
\tag{Demand Inverted}
\end{align*}

This is an isoelastic demand with price-elasticity $\dfrac{1}{1-\alpha}$.

\subsection*{Aggregation and Quality Index}

Define the aggregate quality index:
\[
Q_t := \int_0^N q^{k(s)\frac{\alpha}{1-\alpha}} ds
\]

Then aggregate output simplifies to:
\begin{align*}
\boxed{Y_t = \alpha^{\frac{2\alpha}{1-\alpha}} L\, Q_t}
\tag{Aggregate Output}
\end{align*}

\textbf{Implication:} $Y_t \propto Q_t$, so the growth rate of output equals the growth rate of the quality index on the BGP.

\subsection*{Resource Constraint}
\[
Y_t = C_t + X_t + Z_t
\]
where $X_t = \int_0^N X_t(s)\,ds$ (total intermediates) and $Z_t = \int_0^N Z_t(k(s))\,ds$ (total R\&D).

\clearpage

%==============================================================================
\section{Intermediate Monopolist: Pricing and Profits}
%==============================================================================

\subsection*{Setup}

The successful innovator gains a quality advantage over all previous producers:
\begin{itemize}
    \item \textbf{Old incumbent}: produces quality $q^{k-1}$, marginal cost $= 1$
    \item \textbf{New innovator}: produces quality $q^k$, marginal cost $= 1$
    \item Quality improvement ratio: $q$ (the innovation step)
\end{itemize}

The new innovator solves:
\[
\max_{X_t(s)} \; \pi_t(s) = [P_t(s) - 1]\, X_t(s) \quad \text{subject to demand } X_t(s) = L\left(\frac{\alpha q^{\alpha k(s)}}{P_t(s)}\right)^{\frac{1}{1-\alpha}}
\]

\subsection*{Unconstrained Monopoly Price}

\begin{derivationbox}[title={Derivation: Unconstrained Monopoly Pricing}]
\textbf{Step 1.} Substitute demand into profits:
\[
\pi = (P - 1) \cdot L\left(\frac{\alpha q^{\alpha k}}{P}\right)^{\frac{1}{1-\alpha}}
\]

\textbf{Step 2.} Take FOC with respect to $P$:
\[
\frac{d\pi}{dP} = 0 \implies \frac{1}{P} \cdot \frac{\alpha}{1-\alpha} = \frac{1}{P - 1} \cdot \frac{1}{1-\alpha}
\]

\textbf{Step 3.} Simplify to get the standard markup rule (marginal cost $= 1$, demand elasticity $= 1/(1-\alpha)$):
\[
P^* = \frac{1}{\alpha}
\]
This is the standard markup pricing formula: $P = \text{MC} \times \frac{1/(1-\alpha)}{1/(1-\alpha)-1} = \frac{1}{\alpha}$.
\end{derivationbox}

\subsection*{Limit Pricing Argument}

The new innovator can threaten to undercut the old producer. In quality-adjusted terms, the old producer's effective price is $\frac{P^{\text{old}}}{q^{k-1}}$ whereas the new producer's effective price is $\frac{P^{\text{new}}}{q^k}$. The new innovator captures the entire market if:
\[
\frac{P^{\text{new}}}{q^k} < \frac{P^{\text{old}}}{q^{k-1}} \iff P^{\text{new}} < q \cdot P^{\text{old}}
\]

The old incumbent's lowest possible price is $P^{\text{old}} = 1$ (marginal cost). So the limit price constraint is:
\[
P^{\text{new}} \leq q \quad \text{(limit price)}
\]

\begin{tcolorbox}[colback=orange!10,colframe=orange3,title=The Two Cases: $q \geq 1/\alpha$ vs.\ $q < 1/\alpha$]
\begin{itemize}
    \item \textbf{Case 1: $q \geq 1/\alpha$ (assumed in main model)} \\
          Unconstrained monopoly price $P^* = 1/\alpha \leq q$.\\
          Limit pricing does NOT bind. Innovator charges $P = 1/\alpha$ and is unchallenged.
    \item \textbf{Case 2: $q < 1/\alpha$ (limit pricing binds)} \\
          If innovator charged $1/\alpha > q$, the old incumbent could undercut.\\
          Innovator must charge $P = q$ (Bertrand competition with old producer).\\
          Old and new producer may coexist.
\end{itemize}
\end{tcolorbox}

\subsection*{Monopoly Profits}

\begin{derivationbox}[title={Derivation: Monopoly Profits (Case $q \geq 1/\alpha$)}]
\textbf{Step 1.} With price $P = 1/\alpha$ and MC $= 1$, the markup is:
\[
P - 1 = \frac{1}{\alpha} - 1 = \frac{1-\alpha}{\alpha}
\]

\textbf{Step 2.} Quantity demanded at $P = 1/\alpha$:
\[
X(k(s)) = L\, \alpha^{\frac{2}{1-\alpha}}\, q^{\frac{\alpha k(s)}{1-\alpha}}
\]

\textbf{Step 3.} Profits per sector $s$ at rung $k(s)$:
\[
\pi(k(s)) = \frac{1-\alpha}{\alpha}\, \alpha^{\frac{2}{1-\alpha}}\, L\, q^{\frac{\alpha k(s)}{1-\alpha}} = \pibar\, q^{\frac{\alpha k(s)}{1-\alpha}}
\]
where $\pibar := \frac{1-\alpha}{\alpha}\,\alpha^{\frac{2}{1-\alpha}} L$.

\textbf{Step 4.} In aggregate (using $Y_t = \alpha^{\frac{2\alpha}{1-\alpha}} L Q_t$):
\[
\pi_{\text{total}} = (1-\alpha)\,Y_t / N \quad \text{(per-sector normalized)}
\]
\end{derivationbox}

\begin{resultbox}
\begin{align*}
\boxed{P = \frac{1}{\alpha}, \quad \pi(k(s)) = \pibar\, q^{\frac{\alpha k(s)}{1-\alpha}}, \quad \pibar = \frac{1-\alpha}{\alpha}\,\alpha^{\frac{2}{1-\alpha}} L}
\tag{Monopoly}
\end{align*}
\end{resultbox}

\clearpage

%==============================================================================
\section{Innovation: R\&D, Value of Innovation, and Bellman Equation}
%==============================================================================

\subsection*{Innovation Technology}

A potential entrant invests $Z$ units of output to achieve an innovation at Poisson rate $p$. The innovation probability is specified as:
\[
p(k(s)) = \frac{Z(k(s))}{\xi\, q^{\frac{\alpha(k(s)+1)}{1-\alpha}}}
\]

This means higher quality rungs require more resources per unit of innovation probability (innovations are harder at the frontier).

The innovator's problem is:
\[
\max_{Z \geq 0} \; p(k(s))\, V(k(s)+1) - Z = \frac{Z}{\xi\, q^{\frac{\alpha(k(s)+1)}{1-\alpha}}}\, V(k(s)+1) - Z
\]

\subsection*{Value of Innovation and Bellman Equation}

Once the innovator holds the patent at rung $k$, they earn flow profits $\pi(k)$ until replaced by the next innovator arriving at Poisson rate $p$ (creative destruction). With a constant interest rate $r$:

\begin{derivationbox}[title={Derivation: Value Function (Bellman Equation)}]
\textbf{Step 1.} In continuous time, the value of the patent satisfies the asset pricing (HJB) equation:
\[
r\,V(k) = \pi(k) - p\, V(k) + \dot{V}
\]
The incumbent earns $\pi(k)$ per unit time, loses $V(k)$ when an innovation arrives (rate $p$), and $\dot{V}$ captures capital gains.

\textbf{Step 2.} On the BGP, $V$ grows at the same rate as $\pi$ and $Y$, i.e., $\dot{V}/V = \gamma$. Rearranging:
\[
(r + p - \gamma)\, V(k) = \pi(k) \implies V(k) = \frac{\pi(k)}{r + p - \gamma}
\]

\textbf{Step 3.} Simplification (constant $\pi$, $\dot{V}=0$, constant rate $r$):
\[
V(k) = \frac{\pi(k)}{r + p}
\]

\textbf{Step 4.} Using $\pi(k(s)) = \pibar\, q^{\frac{\alpha k}{1-\alpha}}$:
\[
V(k(s)) = \frac{\pibar\, q^{\frac{\alpha k(s)}{1-\alpha}}}{r + p}
\]
\end{derivationbox}

\begin{resultbox}
\begin{align*}
\boxed{r\,V = \pi - p\,V + \dot{V}} \quad \text{(Bellman / no-arbitrage)}
\tag{HJB} \\[0.5em]
\boxed{V(k) = \frac{\pi(k)}{r + p}} \quad \text{(BGP with } \dot{V}=0\text{)}
\tag{Value}
\end{align*}
\end{resultbox}

\subsection*{Free Entry Condition into R\&D}

With free entry, the expected return to R\&D equals its cost. The FOC of the innovator's problem gives:
\begin{align*}
\frac{V(k(s)+1)}{\xi\, q^{\frac{\alpha(k(s)+1)}{1-\alpha}}} = 1
\tag{Free Entry}
\end{align*}

Since $V(k+1) = \pibar\, q^{\frac{\alpha(k+1)}{1-\alpha}} / (r+p)$, this simplifies to:
\[
\frac{\pibar}{\xi(r+p)} = 1 \implies \boxed{r = \frac{\pibar}{\xi} - p}
\tag{Free Entry Simplified}
\]

This is a key relationship: the interest rate equals the ratio of normalized profits to R\&D costs, minus the innovation rate.

\subsection*{Equilibrium Innovation Rate}

The optimal R\&D investment per sector is:
\[
Z^*(k(s)) = \xi\, q^{\frac{\alpha(k(s)+1)}{1-\alpha}}\, p = q^{\frac{\alpha(k(s)+1)}{1-\alpha}}(\pibar - r\xi)
\]

\textbf{Key insight:} The innovation arrival rate $p = \pibar/\xi - r$ is \emph{identical across all sectors}, regardless of the current quality rung. Higher-rung sectors simply require more resources per unit of $p$.

\clearpage

%==============================================================================
\section{Free Entry, Euler Equation, and BGP Growth Rate}
%==============================================================================

\subsection*{Quality Growth Dynamics}

Each sector independently improves quality by a factor $q$ at Poisson rate $p$. Using the law of large numbers across sectors:
\begin{align*}
\frac{\dot{Q}}{Q} = p\left[q^{\frac{\alpha}{1-\alpha}} - 1\right]
\tag{Quality Growth}
\end{align*}

Since $Y_t \propto Q_t$, the BGP growth rate of output is:
\begin{align*}
\boxed{\gamma = \frac{\dot{Y}}{Y} = \frac{\dot{Q}}{Q} = p\left[q^{\frac{\alpha}{1-\alpha}} - 1\right]}
\tag{BGP Output Growth}
\end{align*}

\subsection*{Household Euler Equation}

Households maximize $\int_0^\infty \frac{c^{1-\sigma}-1}{1-\sigma}e^{-\rho t}dt$, yielding:
\begin{align*}
\frac{\dot{c}}{c} = \frac{1}{\sigma}(r - \rho)
\tag{Euler}
\end{align*}

On the BGP: $\dot{c}/c = \gamma$, so:
\[
r = \rho + \sigma\gamma
\]

\subsection*{Solving for the BGP}

\begin{derivationbox}[title={Derivation: BGP Growth Rate (log utility, $\sigma=1$)}]
\textbf{Step 1.} From free entry: $r = \pibar/\xi - p$.

\textbf{Step 2.} From Euler ($\sigma=1$): $r = \rho + \gamma$.

\textbf{Step 3.} Combine:
\[
\rho + \gamma = \frac{\pibar}{\xi} - p
\]

\textbf{Step 4.} From quality growth: $\gamma = p(q^{\frac{\alpha}{1-\alpha}} - 1)$. Substitute:
\[
\rho + p\left(q^{\frac{\alpha}{1-\alpha}} - 1\right) = \frac{\pibar}{\xi} - p
\implies p\, q^{\frac{\alpha}{1-\alpha}} = \frac{\pibar}{\xi} - \rho
\]
\[
\boxed{p = \left(\frac{\pibar}{\xi} - \rho\right) q^{-\frac{\alpha}{1-\alpha}}}
\]

\textbf{Step 5.} Substitute back for growth rate:
\[
\gamma = p\left(q^{\frac{\alpha}{1-\alpha}} - 1\right) = \left(\frac{\pibar}{\xi} - \rho\right)\left(1 - q^{-\frac{\alpha}{1-\alpha}}\right)
\]
\end{derivationbox}

\begin{resultbox}
\begin{align*}
\boxed{p^* = \left(\frac{\pibar}{\xi} - \rho\right) q^{-\frac{\alpha}{1-\alpha}}}
\tag{Eqm Innovation Rate}\\[0.5em]
\boxed{\gamma^* = \left(\frac{\pibar}{\xi} - \rho\right)\!\left(1 - q^{-\frac{\alpha}{1-\alpha}}\right)}
\tag{BGP Growth}\\[0.5em]
\boxed{r^* = \frac{\pibar}{\xi}\!\left(1 - q^{-\frac{\alpha}{1-\alpha}}\right) + \rho\, q^{-\frac{\alpha}{1-\alpha}}}
\tag{Eqm Interest Rate}
\end{align*}
\end{resultbox}

\subsection*{BGP Existence Condition}

For $p > 0$ and $\gamma > 0$, we need:
\[
\frac{\pibar}{\xi} > \rho \quad \text{(profits must be large relative to discounting)}
\]

\clearpage

%==============================================================================
\section{Market Equilibrium vs Social Planner}
%==============================================================================

\subsection*{Planner's Problem}

The social planner maximizes household welfare subject to resource and technology constraints:
\[
\max \; L\int_0^\infty \frac{c_t^{1-\sigma}-1}{1-\sigma} e^{-\rho t} dt
\]
subject to $Y_t = X_t + Z_t + Lc_t$ and $\dot{Q}_t = \dfrac{Z_t(1-q^{-\frac{\alpha}{1-\alpha}})}{\xi}$.

\subsection*{Planner's Optimal Allocations}

The planner sets price equal to marginal cost ($P^* = 1$ vs market $P = 1/\alpha$), so:
\[
X_t^*(s) = L\,\alpha^{\frac{1}{1-\alpha}}\, q^{\frac{\alpha k(s)}{1-\alpha}} > X_t(s) = L\,\alpha^{\frac{2}{1-\alpha}}\, q^{\frac{\alpha k(s)}{1-\alpha}}
\]

The planner uses \emph{more} intermediates than the market (no monopoly markup distortion). Correspondingly:
\[
Y_t^* = L\,\alpha^{\frac{\alpha}{1-\alpha}} Q_t > Y_t = L\,\alpha^{\frac{2\alpha}{1-\alpha}} Q_t \quad \text{(at given } Q_t\text{)}
\]

The planner's implicit profit from innovation is:
\begin{align*}
\pi^* = \frac{1-\alpha}{\alpha}\,\alpha^{\frac{1}{1-\alpha}} L > \pibar = \frac{1-\alpha}{\alpha}\,\alpha^{\frac{2}{1-\alpha}} L
\tag{Planner Profit}
\end{align*}

Since $\alpha^{\frac{1}{1-\alpha}} > \alpha^{\frac{2}{1-\alpha}}$ for $\alpha \in (0,1)$.

\subsection*{Planner's Growth Rate}

Using the Hamiltonian on the planner's problem yields the planner's Euler equation with implicit interest rate:
\[
r^* = \frac{\pi^*}{\xi}\left(1 - q^{-\frac{\alpha}{1-\alpha}}\right)
\]

And the planner's growth rate:
\begin{align*}
\gamma^* = \frac{1}{\sigma}\left[\frac{\pi^*}{\xi}\left(1 - q^{-\frac{\alpha}{1-\alpha}}\right) - \rho\right]
\tag{Planner Growth}
\end{align*}

\subsection*{Comparing Market and Planner Rates}

Market interest rate:
\[
r = \frac{\pibar}{\xi}\left(1 - q^{-\frac{\alpha}{1-\alpha}}\right) + \rho\, q^{-\frac{\alpha}{1-\alpha}}
\]

\begin{derivationbox}[title={Two Externalities Driving the Wedge}]
\textbf{Externality 1 — Business Stealing (negative for welfare):}\\
An entrant earns the \emph{full} value $V(k+1)$, ignoring that it destroys the incumbent's profits $V(k)$.
In the market, the innovator captures the entire monopoly rent without internalizing the loss imposed on the previous incumbent. This creates an incentive for \emph{excessive} R\&D.
The extra term $\rho\, q^{-\frac{\alpha}{1-\alpha}} > 0$ in $r$ captures this: market interest is too high, so market innovations $p$ are too low given households' discounting --- but entrants innovate \emph{too much} relative to the social value of destroying incumbents.

\textbf{Externality 2 — Appropriability / Knowledge Spillover (positive for welfare):}\\
Innovators capture only monopoly profits $\pibar$, but the social value of innovation includes consumer surplus (the difference between WTP and price). The planner values innovation at $\pi^* > \pibar$ because it sets $P = MC = 1$ rather than $P = 1/\alpha$.
This creates an incentive for \emph{insufficient} R\&D in the market.

\textbf{Net Effect:}
\[
\gamma^* \gtrless \gamma \quad \text{(ambiguous)}
\]
\end{derivationbox}

\begin{resultbox}
\textbf{Business stealing} $\Rightarrow$ tendency toward over-innovation in market.\\
\textbf{Appropriability externality} $\Rightarrow$ tendency toward under-innovation in market.\\
\textbf{Net direction depends on parameters ($q$, $\alpha$, $\rho$).}\\[0.3em]
\textbf{Contrast with Romer:} In expanding varieties models, only the appropriability (positive) externality exists --- no business stealing since old varieties coexist. Result: unambiguous under-innovation.
\end{resultbox}

\subsection*{Policy Implications}

\begin{itemize}
    \item A tax $\tau$ on R\&D raises the cost of entry and reduces $p$ and $\gamma$.
    \item Crucially: \emph{current incumbents favor taxing R\&D} because it protects their market position.
    \item This creates a political economy problem: those with most influence (current monopolists) have the strongest incentive to suppress innovation.
    \item If the economy has too much growth, taxing R\&D may be Pareto-improving. If too little, subsidizing may be.
\end{itemize}

\clearpage

%==============================================================================
\section{Competition and Innovation: Escape Competition Effect}
%==============================================================================

\subsection*{Step-by-Step Innovation with Duopolists}

The Aghion-Howitt (AH) framework can be extended to allow \emph{neck-and-neck} competition between two incumbent firms:

\begin{itemize}
    \item Each product line $j$ is supplied by a duopoly $\{\alpha, \beta\}$ plus a competitive fringe.
    \item Firms' productivities: $A_i = \gamma^{k_i}$, with competitors at most one step apart.
    \item Two types of sectors:
    \begin{itemize}
        \item \textbf{Neck-and-neck} ($\Delta k = 0$): both firms at same rung, fringe one step below.
        \item \textbf{Unlevel} ($\Delta k = 1$): leader one step ahead of follower, fringe two steps below.
    \end{itemize}
\end{itemize}

\subsection*{Bertrand Competition and Payoffs}

Under Bertrand competition:
\begin{itemize}
    \item \textbf{Unlevel sector}: Leader prices at follower's marginal cost $\Rightarrow$ markup $= \gamma$.
    \[
    \pi_1 = \frac{\gamma - 1}{\gamma}, \quad \pi_{-1} = 0
    \]
    \item \textbf{Neck-and-neck sector}: Each firm gets
    \[
    \pi_0 = (1-\Delta)\pi_1
    \]
    where $\Delta \in [0,1]$ is the competition intensity ($\Delta \to 1$: full competition, $\pi_0 \to 0$).
\end{itemize}

\subsection*{Innovation Incentives}

\textbf{Unlevel sector (laggard innovates at rate $n_{-1}$):}
\[
\max_{n_{-1}} (n_{-1} + h)\pi_0 - \frac{\psi n_{-1}^2}{2} \implies n_{-1} = (1-\Delta)\frac{\pi_1}{\psi}
\]
More competition ($\Delta \uparrow$) $\Rightarrow$ $n_{-1} \downarrow$ (Schumpeterian effect).

\textbf{Neck-and-neck sector (one firm innovates at rate $n_0$):}
\[
\max_{n_0} n_0 \pi_1 + (1-n_0)\pi_0 - \frac{\psi n_0^2}{2} \implies n_0 = \frac{\Delta\,\pi_1}{\psi}
\]
More competition ($\Delta \uparrow$) $\Rightarrow$ $n_0 \uparrow$ (\textbf{escape competition effect}).

\subsection*{Inverted-U Relationship}

In steady state, let $\mu_1$ = fraction of unlevel sectors:
\[
\mu_1 = \frac{n_0}{n_0 + n_{-1} + h}
\]

The total innovation flow is:
\[
I = \frac{2\pi_1}{\psi\!\left(1 + \frac{\psi h}{\pi_1}\right)} \,\Delta\!\left(1 - \Delta + \frac{\psi h}{\pi_1}\right)
\]

\begin{derivationbox}[title={Why the Inverted-U Shape?}]
\textbf{Derivative with respect to $\Delta$:}
\[
\frac{\partial I}{\partial \Delta} \propto 1 - 2\Delta + \frac{\psi h}{\pi_1}
\]

\textbf{Low competition ($\Delta$ small):} Most sectors neck-and-neck $\Rightarrow$ escape competition effect dominates $\Rightarrow$ $\partial I / \partial \Delta > 0$.

\textbf{High competition ($\Delta$ near 1):} Most sectors unlevel $\Rightarrow$ Schumpeterian effect dominates $\Rightarrow$ $\partial I / \partial \Delta < 0$.

\textbf{Result:} Non-monotonic (inverted-U) relationship between competition $\Delta$ and innovation $I$.
\end{derivationbox}

\begin{resultbox}
\textbf{Key finding (Aghion et al., 2005):} Competition and innovation have an \emph{inverted-U} relationship:
\begin{itemize}
    \item At low competition: more competition $\Rightarrow$ more innovation (escape competition dominates)
    \item At high competition: more competition $\Rightarrow$ less innovation (Schumpeterian effect dominates)
\end{itemize}
Empirical evidence: UK firm-level data 1973--1994 (citation-weighted patents vs Lerner index).
\end{resultbox}

\clearpage

%==============================================================================
\section{Comparison: Expanding Varieties vs Quality Ladders}
%==============================================================================

\begin{center}
\begin{tabular}{lp{5.5cm}p{5.5cm}}
\toprule
\textbf{Feature} & \textbf{Expanding Varieties (Romer)} & \textbf{Quality Ladders (Aghion-Howitt)} \\
\midrule
Growth mechanism & New goods created; $N$ grows & Quality improvements; $N$ fixed \\
Horizontal / vertical & Horizontal differentiation & Vertical differentiation \\
Creative destruction & No: new goods coexist with old & Yes: new innovator displaces incumbent \\
Business stealing & No (new variety adds to market) & Yes (entrant destroys incumbent's value) \\
Number of varieties & $N$ grows over time & $N$ constant \\
Incumbent fate & Keeps market share, earns forever & Destroyed by next innovator \\
Welfare bias & Under-innovation (appropriability only) & Ambiguous (two opposing externalities) \\
Policy prescription & Subsidize R\&D & Depends on which externality dominates \\
Political economy & Incumbents neutral on innovation & Incumbents actively resist innovation \\
Growth rate depends on & Size of $N$-expansion technology & Quality step $q$, R\&D productivity $1/\xi$ \\
\bottomrule
\end{tabular}
\end{center}

\subsection*{Mathematical Correspondence}

In both models, the BGP growth rate has a similar structure:
\[
\gamma_{\text{Romer}} = \frac{\lambda L - \rho \eta}{\sigma \eta + 1} \quad \text{(from free entry + Euler)}
\]
\[
\gamma_{\text{AH}} = \left(\frac{\pibar}{\xi} - \rho\right)\!\left(1 - q^{-\frac{\alpha}{1-\alpha}}\right) \quad \text{(log utility)}
\]

Both are \emph{scale-effect models}: growth depends on total labor $L$. Both require R\&D profits to cover discounting. Both pin down growth through the free-entry + Euler system.

\subsection*{Key Difference: Conflict of Interest}

\begin{tcolorbox}[colback=orange!10,colframe=orange3,title=The Political Economy Difference]
In Expanding Varieties: no conflict between incumbents and entrants. Everyone benefits from more innovation.\\[0.5em]
In Quality Ladders: innovation \emph{destroys} incumbent profits. Current monopolists have an incentive to lobby against R\&D policies (e.g., support R\&D taxes), even if such policies reduce aggregate welfare. This creates a conflict of interest between incumbents (losers from innovation) and consumers/entrants (gainers from innovation).
\end{tcolorbox}

\clearpage

%==============================================================================
\section{Problem-Solving Templates}
%==============================================================================

\subsection*{(a) Derive Monopoly Price and Profits}

\begin{enumerate}
    \item Write down the final good firm's demand for intermediate $s$:
    \[
    P_t(s) = \alpha L^{1-\alpha} q^{\alpha k(s)} X_t(s)^{\alpha-1}
    \]
    \item The monopolist sets $\max_{X} (P(X)-1) X$. Use the isoelastic demand to get markup $P = \text{MC}/\alpha = 1/\alpha$.
    \item Compute equilibrium quantity: $X^* = L\alpha^{\frac{2}{1-\alpha}} q^{\frac{\alpha k}{1-\alpha}}$.
    \item Compute profits: $\pi = (P-1)X = \frac{1-\alpha}{\alpha} X = \pibar q^{\frac{\alpha k}{1-\alpha}}$.
    \item \textbf{Check} whether $q \geq 1/\alpha$: if yes, limit pricing does not bind; if no, use $P = q$ instead.
\end{enumerate}

\subsection*{(b) Set Up and Solve the Bellman Equation for $V$}

\begin{enumerate}
    \item Write the no-arbitrage HJB: $r\,V = \pi - p\,V + \dot{V}$.
    \item On the BGP, $V$ grows at rate $\gamma$, so $\dot{V} = \gamma V$.
    \item Rearrange: $(r + p - \gamma) V = \pi \implies V = \pi/(r + p - \gamma)$.
    \item Simplified version (often used): $V = \pi/(r + p)$ when $\dot{V} = 0$.
    \item Apply free entry: $V(k+1)/[\xi q^{\frac{\alpha(k+1)}{1-\alpha}}] = 1$, which gives $\pibar/[\xi(r+p)] = 1$.
\end{enumerate}

\subsection*{(c) Apply Free Entry + Euler for BGP Growth}

\begin{enumerate}
    \item \textbf{Free entry:} $r = \pibar/\xi - p$ (from innovator's FOC).
    \item \textbf{Quality growth:} $\gamma = p(q^{\frac{\alpha}{1-\alpha}} - 1)$.
    \item \textbf{Euler equation:} $r = \rho + \sigma\gamma$ (households' optimality).
    \item Substitute Euler into free entry: $\rho + \sigma\gamma = \pibar/\xi - p$.
    \item Substitute quality growth for $\gamma$: solve for $p$ and $\gamma$ simultaneously.
    \item For $\sigma = 1$: $p = (\pibar/\xi - \rho)\, q^{-\frac{\alpha}{1-\alpha}}$ and $\gamma = (\pibar/\xi - \rho)(1 - q^{-\frac{\alpha}{1-\alpha}})$.
\end{enumerate}

\subsection*{(d) Identify and Sign Externalities}

\begin{enumerate}
    \item Solve the planner's problem: Hamiltonian with $Q$ as state, $Z$ (or $C$) as control.
    \item Planner sets $P = MC = 1$ (no markup), yielding $\pi^* > \pibar$.
    \item Compare $r^*$ (planner) vs $r$ (market):
    \begin{itemize}
        \item Term $(\pi^*/\xi - \pibar/\xi)$ reflects appropriability externality (market under-values innovation).
        \item Term $\rho q^{-\frac{\alpha}{1-\alpha}} > 0$ reflects business-stealing (market over-values innovation).
    \end{itemize}
    \item Sign the net externality: $\gamma^* \gtrless \gamma$ depends on whether appropriability or business-stealing dominates.
    \item State policy direction: subsidize R\&D if $\gamma < \gamma^*$; tax if $\gamma > \gamma^*$.
\end{enumerate}

\subsection*{(e) Analyze Competition and Innovation (Escape Competition)}

\begin{enumerate}
    \item Set up payoffs for neck-and-neck ($\pi_0 = (1-\Delta)\pi_1$) and unlevel ($\pi_1$, $0$) sectors.
    \item Solve innovator's optimization in each sector type: get $n_0 = \Delta\pi_1/\psi$ and $n_{-1} = (1-\Delta)\pi_1/\psi$.
    \item Compute steady-state sector composition $\mu_1$ and aggregate innovation flow $I$.
    \item Differentiate $I$ w.r.t. $\Delta$ and sign: $\partial I/\partial\Delta \gtrless 0$ depending on whether $\Delta \lessgtr (1 + \psi h/\pi_1)/2$.
    \item Conclude: inverted-U relationship. Identify which regime the economy is in.
\end{enumerate}

\clearpage

%==============================================================================
\section{Common Pitfalls}
%==============================================================================

\begin{enumerate}
    \item \textbf{Confusing the two $q$ cases for limit pricing}
    \begin{itemize}
        \item If $q \geq 1/\alpha$: monopolist charges $P = 1/\alpha$; limit pricing not binding.
        \item If $q < 1/\alpha$: monopolist charges $P = q$; old and new may coexist.
        \item Always check which case applies before deriving profits.
    \end{itemize}

    \item \textbf{Forgetting that $N$ is fixed}
    \begin{itemize}
        \item Unlike Romer where $N$ grows, in AH the number of sectors $N$ is constant.
        \item Growth comes \emph{entirely} from quality improvements, not variety expansion.
    \end{itemize}

    \item \textbf{Misidentifying the direction of externalities}
    \begin{itemize}
        \item Business stealing: entrant ignores \emph{destroyed} incumbent value $\Rightarrow$ over-innovation tendency.
        \item Appropriability: innovator cannot capture full consumer surplus $\Rightarrow$ under-innovation tendency.
        \item Many students reverse these; remember: ``stealing'' refers to the externality \emph{on the incumbent}.
    \end{itemize}

    \item \textbf{Claiming unambiguous over- or under-innovation}
    \begin{itemize}
        \item In AH, the result is \emph{ambiguous}: either $\gamma > \gamma^*$ or $\gamma < \gamma^*$ is possible.
        \item Only in Romer (expanding varieties) is the direction unambiguous (under-innovation).
    \end{itemize}

    \item \textbf{Errors in the Bellman equation}
    \begin{itemize}
        \item The HJB on the BGP: $r\,V = \pi - p\,V + \dot{V}$. Do NOT omit the $\dot{V}$ term unless you argue it is zero.
        \item On the BGP $\dot{V}/V = \gamma$, so $V = \pi/(r + p - \gamma)$, not $\pi/(r+p)$.
        \item The simplified $V = \pi/(r+p)$ is valid when $\dot{V} = 0$ (constant $V$), which holds if $\pi$ is constant and $r$ is constant.
    \end{itemize}

    \item \textbf{Confusing escape competition with Schumpeterian effects}
    \begin{itemize}
        \item Escape competition applies to \emph{neck-and-neck} sectors: competition increases innovation.
        \item Schumpeterian effect applies to \emph{unlevel} sectors: competition reduces innovation.
        \item Do not say ``competition always reduces innovation'' --- it depends on sector type.
    \end{itemize}

    \item \textbf{Mixing up $p$ (innovation rate) and $\gamma$ (growth rate)}
    \begin{itemize}
        \item $p$ is the Poisson arrival rate of innovations (a frequency).
        \item $\gamma$ is the growth rate of output/quality: $\gamma = p(q^{\frac{\alpha}{1-\alpha}} - 1)$.
        \item They are proportional but not equal.
    \end{itemize}

    \item \textbf{Scale effects}
    \begin{itemize}
        \item Both Romer and AH exhibit scale effects: $\gamma$ depends on total $L$.
        \item This is a known shortcoming; Jones (1995) semi-endogenous growth removes it.
    \end{itemize}
\end{enumerate}

\clearpage

%==============================================================================
\section{Quick Reference Formulas}
%==============================================================================

\subsection*{Production and Demand}
\begin{align*}
Y_t &= L^{1-\alpha}\int_0^N (q^{k(s)} X_t(s))^\alpha ds \tag{Final Good}\\
P_t(s) &= \alpha L^{1-\alpha} q^{\alpha k(s)} X_t(s)^{\alpha-1} \tag{Intermediate Demand}\\
Y_t &= \alpha^{\frac{2\alpha}{1-\alpha}} L\, Q_t, \quad Q_t = \int_0^N q^{k(s)\frac{\alpha}{1-\alpha}}ds \tag{Aggregate Output}
\end{align*}

\subsection*{Monopoly Pricing and Profits}
\begin{align*}
P &= \frac{1}{\alpha} \quad \text{(if } q \geq 1/\alpha\text{)} \tag{Monopoly Price}\\
\pi(k(s)) &= \pibar\, q^{\frac{\alpha k(s)}{1-\alpha}}, \quad \pibar = \frac{1-\alpha}{\alpha}\alpha^{\frac{2}{1-\alpha}}L \tag{Monopoly Profit}
\end{align*}

\subsection*{Innovation and Value}
\begin{align*}
r\,V &= \pi - p\,V + \dot{V} \tag{Bellman / HJB}\\
V(k) &= \frac{\pi(k)}{r + p} \quad (\dot{V}=0 \text{ case}) \tag{Value}\\
\frac{\pibar}{\xi(r+p)} &= 1 \quad \text{(free entry)} \implies r = \frac{\pibar}{\xi} - p \tag{Free Entry}
\end{align*}

\subsection*{BGP (log utility)}
\begin{align*}
\frac{\dot{Q}}{Q} &= p\!\left(q^{\frac{\alpha}{1-\alpha}}-1\right) = \gamma \tag{Quality Growth}\\
p^* &= \left(\frac{\pibar}{\xi} - \rho\right) q^{-\frac{\alpha}{1-\alpha}} \tag{Eqm $p$}\\
\gamma^* &= \left(\frac{\pibar}{\xi} - \rho\right)\!\left(1 - q^{-\frac{\alpha}{1-\alpha}}\right) \tag{BGP Growth}\\
r^* &= \frac{\pibar}{\xi}\!\left(1 - q^{-\frac{\alpha}{1-\alpha}}\right) + \rho\, q^{-\frac{\alpha}{1-\alpha}} \tag{Eqm $r$}
\end{align*}

\subsection*{Planner vs Market}
\begin{align*}
\pi^* &= \frac{1-\alpha}{\alpha}\alpha^{\frac{1}{1-\alpha}}L > \pibar \tag{Planner Profit}\\
r^*_{\text{planner}} &= \frac{\pi^*}{\xi}\!\left(1-q^{-\frac{\alpha}{1-\alpha}}\right) \tag{Planner Rate}\\
\gamma^*_{\text{planner}} &\gtrless \gamma^*_{\text{market}} \quad \text{(ambiguous)} \tag{Welfare}
\end{align*}

\subsection*{Escape Competition (Duopoly Extension)}
\begin{align*}
n_0 &= \frac{\Delta\,\pi_1}{\psi}, \quad n_{-1} = \frac{(1-\Delta)\pi_1}{\psi} \tag{Innovation Rates}\\
I &= \frac{2\pi_1}{\psi\!\left(1+\frac{\psi h}{\pi_1}\right)}\,\Delta\!\left(1-\Delta+\frac{\psi h}{\pi_1}\right) \tag{Aggregate Innovation}\\
\frac{\partial I}{\partial \Delta} &\propto 1 - 2\Delta + \frac{\psi h}{\pi_1} \gtrless 0 \tag{Inverted-U}
\end{align*}

\subsection*{Comparison Table: Key Formulas}
\begin{center}
\begin{tabular}{lcc}
\toprule
\textbf{Object} & \textbf{Romer (EV)} & \textbf{AH (QL)} \\
\midrule
Growth driver & $\dot{N}/N$ & $p(q^{\frac{\alpha}{1-\alpha}}-1)$ \\
Free entry pins & $V_t = \eta w_t$ & $\pibar/[\xi(r+p)] = 1$ \\
Welfare bias & Under-innovation & Ambiguous \\
$\pi$ vs $\pi^*$ & $\pi < \pi^*$ & $\pibar < \pi^*$ \\
Business stealing & None & Yes \\
\bottomrule
\end{tabular}
\end{center}

\end{document}
